\subsection{Expected Values, Covariance, and Correlation}

\hrulefill

\paragraph{Definition}
Let $X$ and $Y$ be jointly distributed random variables with pmf or pdf $P(x,y)$. The \textbf{expected value} of a function $h(X,Y)$ is given by
\[
E[h(X,Y)] = \sum_{x} \sum_{y} h(x,y) P(x,y) \quad \text{(discrete)}
\]
or
\[
E[h(X,Y)] = \int_{-\infty}^{\infty} \int_{-\infty}^{\infty} h(x,y) f(x,y) \, dx \, dy \quad \text{(continuous)}
\] 

Expected values and variances can be used to describe the behavior of a single random variable. If we want to consider joint behavior,
we need to consider how one quantity influence the other?

\subsubsection*{Covariance}
\paragraph{Definition}
The \textbf{covariance} between two random variables $X$ and $Y$ is defined as
\[Cov(X,Y) = E[(X - \mu_X)(Y - \mu_Y)] = E[XY] - \mu_X \mu_Y\]
In practice, the following computational formula is typically used:
\[Cov(X,Y) = E[XY] - E[X]E[Y]\]
The covarience describes the joint behavior of the variables $X$ and $Y$. However, this quantity depends
on the measurement units of $X$ and $Y$.

\subsubsection*{Correlation Coefficient}
The correlation coefficient of $X$ and $Y$ is denoted by $\rho_{X,Y}$ and is defined as
\[\rho_{X,Y} = \frac{Cov(X,Y)}{\sigma_X \sigma_Y}\]
The correlation coefficient is a unitless measure of the linear relationship between $X$ and $Y$.
Its value is always between -1 and 1. A value of 1 indicates a perfect positive linear relationship,
a value of -1 indicates a perfect negative linear relationship, and a value of 0 indicates no linear relationship.

\paragraph{Properties of Covariance and Correlation}
\begin{itemize}
    \item $Cov(X,X) = Var(X)$
    \item $Cov(X,Y) = Cov(Y,X)$
    \item $Cov(aX, bY) = ab \, Cov(X,Y)$
    \item $Cov(X + Y, Z) = Cov(X,Z) + Cov(Y,Z)$
    \item If $X$ and $Y$ are independent, then $Cov(X,Y) = 0$ and $\rho_{X,Y} = 0$
    \item $\rho = \pm 1 \iff \exists a, b \in \mathbb{R} \ni Y = aX+b$ 
\end{itemize}